\documentclass[12pt]{article}

\pagestyle{empty}
\setlength{\topmargin}{0in}
\setlength{\headheight}{0in}
\setlength{\topsep}{0in}
\setlength{\textheight}{9in}
\setlength{\oddsidemargin}{0in}
\setlength{\evensidemargin}{0in}
\setlength{\textwidth}{6.5in}

\usepackage{palatino,graphics,amsmath,tcolorbox,amssymb}

\newcommand{\ds}{\displaystyle}
\newcommand{\vs}[1]{\vspace{#1in}}
\renewcommand{\vss}[1]{\vspace*{#1in}}
\newcommand{\bvec}{{\mathbf b}}
\newcommand{\cvec}{{\mathbf c}}
\newcommand{\dvec}{{\mathbf d}}
\newcommand{\evec}{{\mathbf e}}
\newcommand{\fvec}{{\mathbf f}}
\newcommand{\qvec}{{\mathbf q}}
\newcommand{\uvec}{{\mathbf u}}
\newcommand{\vvec}{{\mathbf v}}
\newcommand{\wvec}{{\mathbf w}}
\newcommand{\xvec}{{\mathbf x}}
\newcommand{\yvec}{{\mathbf y}}
\newcommand{\zvec}{{\mathbf y}}
\newcommand{\zerovec}{{\mathbf 0}}
\newcommand{\s}{\phantom{-}}
\newcommand{\real}{{\mathbb R}}
\newcommand{\twovec}[2]{\left[\begin{array}{r}#1 \\ #2
    \end{array}\right]}
\newcommand{\ctwovec}[2]{\left[\begin{array}{c}#1 \\ #2
   \end{array}\right]}
\newcommand{\threevec}[3]{\left[\begin{array}{r}#1 \\ #2 \\ #3
  \end{array}\right]}
\newcommand{\cthreevec}[3]{\left[\begin{array}{c}#1 \\ #2 \\ #3
    \end{array}\right]}
\newcommand{\fourvec}[4]{\left[\begin{array}{r}#1 \\ #2 \\ #3 \\ #4
    \end{array}\right]}
\newcommand{\cfourvec}[4]{\left[\begin{array}{c}#1 \\ #2 \\ #3 \\ #4
    \end{array}\right]}
\newcommand{\mattwo}[4]{\left[\begin{array}{rr}#1 \amp #2 \\ #3 \amp #4 \\ \end{array}\right]}
\newcommand{\laspan}{\text{Span}}

\begin{document}

\noindent
{\bf Mathematics 204} \\ 
{\bf Linear independence (Section 2.4)}

\bigskip
\begin{enumerate}
\item Consider the matrices 
  $$
  A = 
  \left[
    \begin{array}{cccc}
      3 & 0 & -1 & 1 \\
      1 & -1 & 3 & 7 \\
      3 & -2 & 1 & 5 \\
      -1 & 2 & 2 & 3 \\
    \end{array}
  \right],\hspace*{24pt}
  B = 
  \left[
    \begin{array}{cccc}
      3 & 0 & -1 & 4 \\
      1 & -1 & 3 & -1 \\
      3 & -2 & 1 & 3 \\
      -1 & 2 & 2 & 1 \\
    \end{array}
  \right].
  $$

  Are the columns of $A$ linearly independent?  Are the columns
  of $B$ linearly independent?   Explain your thinking.
  
  \vs{1.5}
  For the matrix above whose columns are linearly dependent,
  express one of the vectors as a linear combination of the others.

  \vs{1.25}
\item Are the three vectors below linearly independent or dependent?
  Explain your thinking.

  \includegraphics[scale=0.7]{figures/lin-dep.pdf}

% \item  Suppose that $\vvec_1$, $\vvec_2$, $\vvec_3$, and $\vvec_4$ are
%   vectors in $\real^8$ and that $\vvec_2=\zerovec$.  What can you say
%   about the linear dependence/independence of this set of
%   four vectors in $\real^8$?
  
%   \vs{1}
%   If the vectors are linearly dependent, express one of them as a
%   linear combination of the others.
  
\vs{1}

\item If $\vvec_1,\vvec_2,\vvec_3$ is a set of vectors in
  $\real^5$, what is the shape of the matrix
  $\begin{bmatrix}\vvec_1 & \vvec_2 & \vvec_3 \end{bmatrix}$ (that is,
  how many rows and columns)?  If the set of vectors is linearly
  independent, what can you say about the pivot positions?

  \vs{1.5}
  
  Explain why seven vectors in $\real^5$ cannot be linearly
  independent.

  \vs{1.25}
  What is the largest number of vectors in $\real^5$ that can be
  linearly independent?

  \vs{1.25}
  If you have three vectors in $\real^5$, can you guarantee that they
  are linearly independent?

  \vs{1}
\item If $A$ is a matrix, we call the equation $A\xvec=\zerovec$ a
  {\em homogeneous equation} since the right hand side is the zero
  vector.  Explain why this equation is always consistent.

  \vs{1}
  \newpage
  Suppose that 
  $
  A = 
  \begin{bmatrix}
    3 & 0 & -1 & 4 \\
    1 & -1 & 3 & -3 \\
    3 & -2 & 1 & 0 \\
    -1 & 2 & 2 & -1
  \end{bmatrix}
  $.  Write the solutions to the equation $A\xvec=\zerovec$ in
  vector-parametric form.

  \vs{2}
  If the columns of $A$ are $\vvec_1$, $\vvec_2$, $\vvec_3$, and
  $\vvec_4$, use your solution to this homogeneous equation to find
  one set of weights such that
  $$
  c_1\vvec_1+c_2\vvec_2+c_3\vvec_3+c_4\vvec_4 = \zerovec
  $$
  where at least one of the weights is nonzero.

  \vs{1.25}
  Use this expression to rewrite $\vvec_4$ as a linear combination of
  the others.

  \vs{1.25}
  Is the set of vectors $\vvec_1$, $\vvec_2$, $\vvec_3$, $\vvec_4$
  linearly dependent or independent?

  \vs{1.25}
  \newpage
\item Suppose that $\wvec_1$, $\wvec_2$, and $\wvec_3$ are the columns
  of a matrix $B$ and that
  $$
  \wvec_3 = 2\wvec_1 - 3\wvec_2.
  $$
  Is this set of vectors linearly dependent or independent?

  \vs{1}
  Explain why we can find a set of weights $c_1$, $c_2$, and $c_3$,
  not all of which are zero, for which
  $$
  c_1\wvec_1+c_2\wvec_2 + c_3\wvec_3 = \zerovec.
  $$

  \vs{1}
  Explain why there are infinitely many solutions to the homogeneous
  equation $B\xvec=\zerovec$.

  \vs{1}
  What does this say about the pivot positions of $B$?

  \vs{1}

\item Suppose that $\vvec_1, \vvec_2, \ldots, \vvec_n$ are vectors in
  $\real^{24}$ that are linearly independent and that span $\real^{24}$.
  What can you say about the pivot positions of the matrix $A$ whose
  columns are $\vvec_1, \vvec_2, \ldots, \vvec_n$?

  \vs{1}
  How many vectors are there in this set?

  \vs{1}
  Suppose that $\bvec$ is a vector in $\real^{24}$.  Explain how you
  know that $\bvec$ can be written as a linear combination of
  $\vvec_1,\vvec_2, \ldots,\vvec_n$.

  \vs{1}
  Suppose that $\bvec$ is a vector in $\real^{24}$.  Explain how you
  know that $\bvec$ can be written as a linear combination of
  $\vvec_1,\vvec_2, \ldots,\vvec_n$ in exactly one way.

  \vs{1}
%  What can you say about the solution to the homogeneous equation
%  $A\xvec = \zerovec$?
%  
%  \vs{1}
\item Suppose that $\vvec_1,\vvec_2,\ldots,\vvec_n$ are a set of
  vectors and that $\vvec_3$ is a scalar multiple of $\vvec_7$.  What
  can you say about the linear independence/dependence of this set of
  vectors?

  \vs{1}
  If $\vvec_1,\vvec_2,\ldots,\vvec_n$ is a linearly independent set of
  vectors, is it necessarily true that one vector is a scalar multiple
  of another one?

  \vs{1}
\item Suppose that $A$ is a matrix such that the equation
  $A\xvec=\bvec$ has exactly one solution for every vector $\bvec$.
  What can you say about the linear dependence/independence of the
  columns of $A$?

  \vs{1}
  What can you say about the span of the columns of $A$?

  \vs{1}
  \begin{tcolorbox}
    Make sure you can answer the following questions:
    \begin{itemize}
    \item What does it mean for a set of vectors to be linearly
      dependent or independent?
    \item How do the pivots of a matrix tell us about the linear
      independence/dependence of its columns?
    \item What is the largest possible number of vectors in a linearly
      independent set in $\real^n$?
    \item How do the solutions to the equation $A\xvec=\zerovec$ give
      information about the linear independence/dependence of the
      columns of $A$?
    \end{itemize}
  \end{tcolorbox}
  
  

\end{enumerate}
\end{document}

  \newpage
\item Suppose you have an $m\times n$ matrix $A$ whose columns span
  $\real^m$.  What can you say about the pivot positions of $A$?

  \vs{1}
  If $\bvec$ is a vector in $\real^m$, what can you say about the
  solutions to $A\xvec = \bvec$?

  \vs{1}
  What can you say about the number of columns of $A$?

  \vs{1}
\item Suppose you have an $m\times n$ matrix $A$ whose columns are
  linearly independent.  What can you say about the pivot positions of
  $A$? 

  \vs{1}
  What can you say about the
  solutions to $A\xvec = \zerovec$?

  \vs{1}
  What can you say about the number of columns of $A$?

  \vs{1}
  
\end{enumerate}
  \vs{1}
\end{document}

